% !TEX program = xelatex
\documentclass[11pt,a4paper]{article}
\usepackage[UTF8]{ctex}
\usepackage{amsmath,amssymb,bm}
\usepackage[margin=2.5cm]{geometry}
\usepackage{booktabs}
\usepackage{enumitem}

\title{Q08 解答}
\author{}
\date{}

\begin{document}
\maketitle


\section{(a)}

设量子态 $|\Psi\rangle \in \mathcal{H}_{\text{imp}} \otimes \mathcal{H}_{\text{env}}$，在基 $\{|a_i\rangle\}$ 和 $\{|b_j\rangle\}$ 下展开：
\begin{equation}
  |\Psi\rangle = \sum_{i=1}^{d_{\text{imp}}} \sum_{j=1}^{d_{\text{env}}} c_{ij}\, |a_i\rangle \otimes |b_j\rangle
\end{equation}

系数矩阵 $C = (c_{ij})$ 的维度为 $d_{\text{imp}} \times d_{\text{env}}$，其中 $d_{\text{imp}} = \dim\mathcal{H}_{\text{imp}}$，$d_{\text{env}} = \dim\mathcal{H}_{\text{env}}$。

对 $C$ 做奇异值分解 $C = U\Sigma V^\dagger$，代入并换基：
\begin{equation}
  |\Psi\rangle = \sum_{k=1}^{r} \lambda_k\, |u_k\rangle \otimes |v_k\rangle
\end{equation}

其中 $r = \text{rank}(C)$ 为 Schmidt 秩，$\lambda_k \geq 0$ 为 Schmidt 系数，$\{|u_k\rangle\}$ 和 $\{|v_k\rangle\}$ 各自正交归一。

\begin{equation}
  \boxed{\;r \leq \min(\dim\mathcal{H}_{\text{imp}},\; \dim\mathcal{H}_{\text{env}})\;}
\end{equation}

$C$ 最多有 $\min(d_{\text{imp}}, d_{\text{env}})$ 个非零奇异值。

HF 态的 1-RDM $\gamma$ 是投影算符（幂等）：$\gamma^2 = \gamma$。

将 $\gamma$ 按碎片 (imp) 和环境 (env) 分块：
\begin{equation}
  \gamma = \begin{pmatrix} \gamma_{\text{II}} & \gamma_{\text{IE}} \\ \gamma_{\text{EI}} & \gamma_{\text{EE}} \end{pmatrix}
\end{equation}

幂等条件 $\gamma^2 = \gamma$ 的左下块给出：
\begin{equation}
  \gamma_{\text{EI}}\,\gamma_{\text{II}} + \gamma_{\text{EE}}\,\gamma_{\text{EI}} = \gamma_{\text{EI}}
\end{equation}

移项：
\begin{equation}
  (I - \gamma_{\text{EE}})\,\gamma_{\text{EI}} = \gamma_{\text{EI}}\,\gamma_{\text{II}}
  \label{eq:idem}
\end{equation}

对两边取秩。右边的秩：
\begin{equation}
  \text{rank}(\gamma_{\text{EI}}\,\gamma_{\text{II}}) \leq \min\big(\text{rank}(\gamma_{\text{EI}}),\; \text{rank}(\gamma_{\text{II}})\big) \leq \text{rank}(\gamma_{\text{II}}) \leq n_{\text{imp}}
\end{equation}

左边的秩：当 $I - \gamma_{\text{EE}}$ 可逆时，$\text{rank}((I-\gamma_{\text{EE}})\gamma_{\text{EI}}) = \text{rank}(\gamma_{\text{EI}})$。即使不可逆，秩只减不增，故 $\text{rank}(\gamma_{\text{EI}}) \geq \text{rank}(\text{左边}) = \text{rank}(\text{右边})$。

结合等式 \eqref{eq:idem}：
\begin{equation}
  \text{rank}(\gamma_{\text{EI}}) = \text{rank}(\gamma_{\text{EI}}\,\gamma_{\text{II}}) \leq \text{rank}(\gamma_{\text{II}}) \leq n_{\text{imp}}
\end{equation}

$\gamma_{\text{EI}}$ 的秩 = 碎片-环境独立耦合方向数 = 浴轨道数 $n_{\text{bath}}$。

\begin{equation}
  \boxed{\;n_{\text{bath}} = \text{rank}(\gamma_{\text{EI}}) \leq n_{\text{imp}}\;}
\end{equation}

对于闭壳层 HF，通常取等号：$n_{\text{bath}} = n_{\text{imp}}$。

\section{(b)}

对角化 $\gamma_{\text{II}}$ 得碎片自然轨道 $\{v_k\}$ 和占据数 $\{\lambda_k\}$：
\begin{equation}
  \gamma_{\text{II}}\, v_k = \lambda_k\, v_k, \qquad k = 1, \ldots, n_{\text{imp}}
\end{equation}

对每个 $v_k$（升维为 $\tilde{v}_k$），构造耦合向量：
\begin{equation}
  w_k = (I - P_{\text{imp}})\,\gamma\,\tilde{v}_k
\end{equation}

归一化得浴轨道 $b_k = w_k / |w_k|$。

\begin{align}
  |w_k|^2 &= \tilde{v}_k^\dagger\, \gamma\, (I - P_{\text{imp}})\, \gamma\, \tilde{v}_k \\
  &= \tilde{v}_k^\dagger\, \gamma^2\, \tilde{v}_k - \tilde{v}_k^\dagger\, \gamma\, P_{\text{imp}}\, \gamma\, \tilde{v}_k \\
  &= \underbrace{\tilde{v}_k^\dagger\, \gamma\, \tilde{v}_k}_{= \lambda_k\ (\text{特征方程})} - \underbrace{|\gamma_{\text{II}} v_k|^2}_{= \lambda_k^2\ (\text{特征方程})} \qquad (\text{用了 } \gamma^2 = \gamma) \\
  &= \lambda_k - \lambda_k^2 = \lambda_k(1 - \lambda_k)
\end{align}

\begin{itemize}
  \item 耦合强度：$c_k = |w_k| = \sqrt{\lambda_k(1-\lambda_k)}$（第二步用 $\gamma^2 = \gamma$ 算出）
  \item 浴占据数：$\mu_k = 1 - \lambda_k$（2$\times$2 块幂等条件给出）
\end{itemize}

在基 $\{v_1, \ldots, v_{n_{\text{imp}}}, b_1, \ldots, b_{n_{\text{imp}}}\}$ 下，全分子 $\gamma$ 限制到嵌入空间的形式是 $n_{\text{imp}}$ 个\textbf{独立的 $2 \times 2$ 块}：
\begin{equation}
  \gamma\big|_{\text{embed}} = \bigoplus_{k=1}^{n_{\text{imp}}} \begin{pmatrix} \lambda_k & c_k \\ c_k & \mu_k \end{pmatrix}
  = \bigoplus_{k=1}^{n_{\text{imp}}} \begin{pmatrix} \lambda_k & \sqrt{\lambda_k(1-\lambda_k)} \\ \sqrt{\lambda_k(1-\lambda_k)} & 1-\lambda_k \end{pmatrix}
\end{equation}

``独立''是因为：碎片自然轨道之间无耦合（$\gamma_{\text{II}}$ 已对角化）；浴轨道之间也无耦合（浴由 $\gamma_{\text{EI}}$ 的列空间构造，不同 $k$ 对应不同列方向，在 $\gamma^2 = \gamma$ 下正交）。

每个 $2 \times 2$ 块的\textbf{左上角}就是 $\lambda_k$。因此嵌入空间密度矩阵的碎片块为：
\begin{equation}
  \big(\gamma\big|_{\text{embed}}\big)_{\text{II}} = \text{diag}(\lambda_1, \lambda_2, \ldots, \lambda_{n_{\text{imp}}}) = \gamma_{\text{II}}^{\text{full}}
\end{equation}

最后一步用了：$\gamma_{\text{II}}^{\text{full}}$ 在自然轨道基下也是 $\text{diag}(\lambda_1, \ldots, \lambda_{n_{\text{imp}}})$（第一步对角化的定义）。

\begin{enumerate}
  \item \textbf{$\gamma\big|_{\text{embed}}$ 幂等}：每个 $2 \times 2$ 块满足 $\begin{pmatrix} \lambda & c \\ c & 1{-}\lambda \end{pmatrix}^2 = \begin{pmatrix} \lambda & c \\ c & 1{-}\lambda \end{pmatrix}$（直接验证：$\lambda^2 + c^2 = \lambda^2 + \lambda(1{-}\lambda) = \lambda$ ✓），故整体幂等 → 对应 Slater 行列式。

  \item \textbf{嵌入空间是 $\gamma$-不变的}：$\gamma$ 作用在 $v_k$ 上 $\to \lambda_k v_k + c_k b_k$（碎片 ∪ 浴，不漏出）；$\gamma$ 作用在 $b_k$ 上 $\to c_k v_k + \mu_k b_k$（同样不漏出，用 $(I - \gamma_{\text{EE}})\gamma_{\text{EI}} = \gamma_{\text{EI}}\gamma_{\text{II}}$ 推出）。所以 $\gamma$ 把嵌入空间映射到自身。

  \item \textbf{Fock 方程保持}：全分子 HF 满足 $[F, \gamma] = 0$。因为嵌入空间 $\gamma$-不变，$F$ 和 $\gamma$ 都可干净地限制到嵌入空间，对易关系保持：$[F\big|_{\text{embed}}, \gamma\big|_{\text{embed}}] = 0$。这就是嵌入空间的 HF 方程。
\end{enumerate}

因此 $\gamma\big|_{\text{embed}}$ 是 $H_{\text{embed}}$ 的 HF 密度矩阵，其碎片块：
\begin{equation}
  \boxed{\;\gamma_{\text{II}}^{\text{embed}} = \text{diag}(\lambda_1, \ldots, \lambda_{n_{\text{imp}}}) = \gamma_{\text{II}}^{\text{full}}\;}
\end{equation}

\section{(c) }

\subsection{为什么 $\gamma^{\text{corr}} \neq (\gamma^{\text{corr}})^2$}

HF 态是单个 Slater 行列式，1-RDM 是投影算符：$\gamma^{\text{HF}}$ 的特征值为 0 或 1，满足 $\gamma^2 = \gamma$。

CCSD/SQD 给出的是\textbf{关联态}，不再是单个 Slater 行列式。关联态的 1-RDM $\gamma^{\text{corr}}$ 的特征值（自然轨道占据数）在 0 和 1 之间取任意值，例如 $n_p = 0.83, 0.12, 0.97$ 等。

验证：若 $\gamma$ 的特征值为 $n_p$，则 $\gamma^2$ 的特征值为 $n_p^2$。$\gamma^2 = \gamma$ 要求 $n_p^2 = n_p$，即 $n_p \in \{0, 1\}$。关联态 $n_p \notin \{0, 1\}$，故：
\begin{equation}
  \boxed{\;(\gamma^{\text{corr}})^2 \neq \gamma^{\text{corr}}\;}
\end{equation}

\subsection{为什么最小浴不再精确}

(b) 的证明\textbf{完全依赖} $\gamma^2 = \gamma$。具体来说：

\begin{enumerate}
  \item 浴数上界 $n_{\text{bath}} \leq n_{\text{imp}}$ 来自 $\gamma^2 = \gamma$ 给出的秩不等式 \eqref{eq:idem}。对关联态，$(\gamma^{\text{corr}})^2 \neq \gamma^{\text{corr}}$，该不等式不成立，Schmidt 秩可能超过 $n_{\text{imp}}$。

  \item 浴轨道属性（占据数 $\mu_k = 1-\lambda_k$，耦合 $c_k = \sqrt{\lambda_k(1-\lambda_k)}$）由 $\gamma^2 = \gamma$ 的 2$\times$2 块推导。对关联态，2$\times$2 块不再满足 $\gamma^2 = \gamma$，这些公式失效。

  \item 配对结构（每个碎片自然轨道恰好配一个浴轨道）依赖 $\gamma^2 = \gamma$。关联态的纠缠结构更复杂——一个碎片自然轨道可能和多个环境方向耦合。
\end{enumerate}

\textbf{结论}：用 HF 构造的 $n_{\text{imp}}$ 个浴轨道，对关联态来说可能\textbf{不够}——丢失了部分碎片-环境纠缠信息。

\subsection{MP2 浴补充了什么关联}

HF 浴编码的是\textbf{平均场级别}的碎片-环境纠缠（静态关联）。缺失的是\textbf{动力学关联}——电子之间的瞬时涨落耦合。

MP2（二阶微扰理论）通过微扰修正计算关联态的 1-RDM 修正 $\delta\gamma$：
\begin{equation}
  \gamma^{\text{corr}} \approx \gamma^{\text{HF}} + \delta\gamma^{\text{MP2}}
\end{equation}

MP2 浴自然轨道（BNO）从 $\delta\gamma^{\text{MP2}}$ 的环境部分提取。BNO 的占据数大致呈指数衰减——前几个 BNO 携带大部分动力学关联信息。

\textbf{MP2 浴补充的内容}：
\begin{itemize}
  \item HF 浴只捕获了 $\gamma^{\text{HF}}$ 的碎片-环境耦合（静态部分）
  \item MP2 浴额外捕获 $\delta\gamma^{\text{MP2}}$ 的碎片-环境耦合（动力学部分）
  \item 截断阈值 $\eta$ 控制 BNO 数量：$\eta$ 越小，浴越大，精度越高但代价越大
\end{itemize}

\section*{DMET 化学势 $\mu$ 与 one-shot vs 自洽}

\subsection*{$\mu$ 的确定}

$\mu$ 是拉格朗日乘子，约束碎片电子数：
$$N_{\text{imp}}(\mu) = \text{Tr}_{\text{imp}}\big[\gamma^{\text{corr}}(\mu)\big] = N_{\text{imp}}^{\text{target}}$$

构造嵌入哈密顿量 $H_{\text{imp}}(\mu) = H_{\text{embed}} + \mu\,\hat{N}_{\text{imp}}$。$\hat{N}_{\text{imp}}$ 在计算基下是对角矩阵，故 $\mu$ 只改变 CI 矩阵对角元。$N_{\text{imp}}(\mu)$ 随 $\mu$ 单调递减，用牛顿法搜索。

\subsection*{One-shot DMET}

用 $\gamma^{\text{HF}}$ 构造浴（一次性），每碎片求解一次，$\mu = 0$ 不迭代。极快但碎片电子数可能不匹配，精度受限。

\subsection*{自洽 DMET}

对每碎片迭代：构造 $H_{\text{imp}}(\mu)$ $\to$ 求解 $\to$ 检查 $N_{\text{imp}}$ $\to$ 更新 $\mu$，3--5 轮收敛。精度更高但慢 3--5 倍。可选浴自洽（用 $\gamma^{\text{corr}}$ 重做浴），计算量翻数倍。

\begin{center}
\begin{tabular}{lcc}
  \toprule
  & One-shot & 自洽 \\
  \midrule
  求解器调用 & $F$（每碎片 1 次） & $\sim 4F$ \\
  $\mu$ 匹配 & 不保证 & 保证 \\
  精度 & 较低 & 较高 \\
  \bottomrule
\end{tabular}
\end{center}

\end{document}
